\section{AIME I 2010}
        \begin{enumerate}
        	\item (\textit{AIME I 2010}) Maya lists all the positive divisors of $2010^2$. She then randomly selects two distinct divisors from this list. Let $p$ be the probability that exactly one of the selected divisors is a perfect square. The probability $p$ can be expressed in the form $\frac {m}{n}$, where $m$ and $n$ are relatively prime positive integers. Find $m + n$.


 
 \textbf{Answer: }107



\textbf{Solution 1}

$2010^2 = 2^2\cdot3^2\cdot5^2\cdot67^2$. Thus there are $(2+1)^4$ divisors, $(1+1)^4$ of which are squares (the exponent of each prime factor must either be $0$ or $2$). Therefore the probability is
 \[\frac {2\cdot2^4\cdot(3^4 - 2^4)}{3^4(3^4 - 1)} = \frac {26}{81} \Longrightarrow 26+ 81 = \boxed{107}.\]




\textbf{Solution 2 (Using a Bit More Counting)}

The prime factorization of $2010^2$ is $67^2\cdot3^2\cdot2^2\cdot5^2$. Therefore, the number of divisors of $2010^2$ is $3^4$ or $81$, $16$ of which are perfect squares. The number of ways we can choose $1$ perfect square from the two distinct divisors is $\binom{16}{1}\binom{81-16}{1}$. The total number of ways to pick two divisors is $\binom{81}{2}$

 Thus, the probability is 
 \[\frac {\binom{16}{1}\binom{81-16}{1}}{\binom{81}{2}} = \frac {16\cdot65}{81\cdot40} = \frac {26}{81} \Longrightarrow 26+ 81 = \boxed{107}.\]




\textbf{Solution 3 (Using a Bit Bit more counting)}

The prime factorization of $2010^2$ is $67^2\cdot3^2\cdot2^2\cdot5^2$. So, there are $3^4 = 81$ divisors. Now, consider the subset $\{2,3,5,67 \}$. Take a subset of this set. Then, square and multiply the elements. We have a perfect square that is a divisor of $2010^2$. From that information, we include the empty set as $\{ \} = 1$ is a viable option. That is 16 possible squares.

 So, the probability (as in solution 1) is $\frac{16}{81} \cdot \frac{80-15}{80} = \frac{16}{81} \cdot \frac{13}{16} = \frac{13}{81}$. Times two due to the number of permutations, we have $p = \frac{26}{81} \implies \boxed{107}$.

 


	\item (\textit{AIME I 2010}) Find the remainder when $9 \times 99 \times 999 \times \cdots \times \underbrace{99\cdots9}_{\text{999 9's}}$ is divided by $1000$.


 
 \textbf{Answer: }109



\textbf{Solution}

Note that $999\equiv 9999\equiv\dots \equiv\underbrace{99\cdots9}_{\text{999 9's}}\equiv -1 \pmod{1000}$ (see \href{https://artofproblemsolving.com/wiki/index.php?title=Modular\_arithmetic}{modular arithmetic}). That is a total of $999 - 3 + 1 = 997$ integers, so all those integers multiplied out are congruent to $- 1\pmod{1000}$. Thus, the entire expression is congruent to $- 1\times9\times99 = - 891\equiv\boxed{109}\pmod{1000}$.

 



\textbf{Solution 2}

The expression also equals $(10-1)(100-1)\dots({10^{999}}-1)$. To find its modular 1,000, remove all terms from 1,000 and after. Then the expression becomes $(10-1)(100-1)(-1) \pmod{1000} \equiv -891 \pmod{1000} \equiv \boxed{109}\pmod{1000}$

 By maxamc

 


	\item (\textit{AIME I 2010}) Suppose that $y = \frac34x$ and $x^y = y^x$. The quantity $x + y$ can be expressed as a rational number $\frac {r}{s}$, where $r$ and $s$ are relatively prime positive integers. Find $r + s$.


 
 \textbf{Answer: }529



\textbf{Solution 1}

Substitute $y = \frac34x$ into $x^y = y^x$ and solve.
 \[x^{\frac34x} = \left(\frac34x\right)^x\]
 \[x^{\frac34x} = \left(\frac34\right)^x \cdot x^x\]
 \[x^{-\frac14x} = \left(\frac34\right)^x\]
 \[x^{-\frac14} = \frac34\]
 \[x = \frac{256}{81}\]
 \[y = \frac34x = \frac{192}{81}\]
 \[x + y = \frac{448}{81}\]
 \[448 + 81 = \boxed{529}\]




\textbf{Solution 2}

We solve in general using $c$ instead of $3/4$. Substituting $y = cx$, we have:

 
 \[x^{cx} = (cx)^x \Longrightarrow (x^x)^c = c^x\cdot x^x\]

 Dividing by $x^x$, we get $(x^x)^{c - 1} = c^x$.

 Taking the $x$th root, $x^{c - 1} = c$, or $x = c^{1/(c - 1)}$.

 In the case $c = \frac34$, $x = \Bigg(\frac34\Bigg)^{ - 4} = \Bigg(\frac43\Bigg)^4 = \frac {256}{81}$, $y = \frac {64}{27}$, $x + y = \frac {256 + 192}{81} = \frac {448}{81}$, yielding an answer of $448 + 81 = \boxed{529}$.

 



\textbf{Solution 3}

Taking the logarithm base $x$ of both sides, we arrive with:

 
 \[y = \log_x y^x \Longrightarrow \frac{y}{x} = \log_{x} y = \log_x \frac{3}{4}x = \frac{3}{4}\]

 Where the last two simplifications were made since $y = \frac{3}{4}x$. Then,

 
 \[x^{\frac{3}{4}} = \frac{3}{4}x \Longrightarrow x^{\frac{1}{4}} = \frac{4}{3} \Longrightarrow x = \left(\frac{4}{3}\right)^4\]

 Then, $y = \left(\frac{4}{3}\right)^3$, and thus:

  
 \[x+y = \left(\frac{4}{3}\right)^3 \left(\frac{4}{3} + 1 \right) = \frac{448}{81} \Longrightarrow 448 + 81 = \boxed{529}\]




\textbf{Solution 4 (another version of Solution 3)}

Taking the logarithm base $x$ of both sides, we arrive with:
 \[y = \log_x y^x \Longrightarrow \frac{y}{x} = \log_{x} y = \log_x \left(\frac{3}{4}x\right) = \frac{3}{4}\]Now we proceed by the logarithm rule $\log(ab)=\log a + \log b$. The equation becomes:
 \[\log_x \frac{3}{4} + \log_x x = \frac{3}{4}\]
 \[\Longleftrightarrow \log_x \frac{3}{4} + 1 = \frac{3}{4}\]
 \[\Longleftrightarrow \log_x \frac{3}{4} = -\frac{1}{4}\]
 \[\Longleftrightarrow x^{-\frac{1}{4}} = \frac{3}{4}\]
 \[\Longleftrightarrow \frac{1}{x^{\frac{1}{4}}} = \frac{3}{4}\]
 \[\Longleftrightarrow x^{\frac{1}{4}} = \frac{4}{3}\]
 \[\Longleftrightarrow \sqrt[4]{x} = \frac{4}{3}\]
 \[\Longleftrightarrow x = \left(\frac{4}{3}\right)^4=\frac{256}{81}\]Then find $y$ as in solution 3, and we get $\boxed{529}$.

 


	\item (\textit{AIME I 2010}) Jackie and Phil have two fair coins and a third coin that comes up heads with probability $\frac47$. Jackie flips the three coins, and then Phil flips the three coins. Let $\frac {m}{n}$ be the probability that Jackie gets the same number of heads as Phil, where $m$ and $n$ are relatively prime positive integers. Find $m + n$.


 
 \textbf{Answer: }515



\textbf{Solution 1}

This can be solved quickly and easily with \href{https://artofproblemsolving.com/wiki/index.php?title=Generating\_functions}{generating functions}.

 Let $x^n$ represent flipping $n$ heads.

 The generating functions for these coins are $(1+x)$,$(1+x)$,and $(3+4x)$ in order.

 The product is $3+10x+11x^2+4x^3$. ($ax^n$ means there are $a$ ways to get $n$ heads, eg there are $10$ ways to get $1$ head, and therefore $2$ tails, here.)

 The sum of the coefficients squared (total number of possible outcomes, squared because the event is occurring twice) is $(4 + 11 + 10 + 3)^2 = 28^2 = 784$ and the sum of the squares of each coefficient (the sum of the number of ways that each coefficient can be chosen by the two people) is $4^2 + 11^2 + 10^2 + 3^2=246$.The probability is then $\frac{4^2 + 11^2 + 10^2 + 3^2}{28^2}  = \frac{246}{784} = \frac{123}{392}$.(Notice the relationship between the addends of the numerator here and the cases in the following solution.)

 $123 + 392 = \boxed{515}$




\textbf{Solution 2}

We perform \href{https://artofproblemsolving.com/wiki/index.php?title=Casework}{casework} based upon the number of heads that are flipped. 

 
            \begin{itemize}
        \item \textbf{Case 1}: No heads.
            \end{itemize}
            
            \begin{enumerate}[label=(\roman*)]
                \item The only possibility is TTT (the third coin being the unfair coin). The probability for this to happen to Jackie is $\frac {1}{2} \cdot \frac {1}{2} \cdot \frac {3}{7} = \frac {3}{28}$ Thus the probability for this to happen to both players is $\left(\frac {3}{28}\right)^2 = \frac {9}{784}$
            \end{enumerate}
            
            \begin{itemize}
        \item \textbf{Case 2}: One head.
            \end{itemize}
            
            \begin{enumerate}[label=(\roman*)]
                \item We can have either HTT, THT, or TTH. The first two happen to Jackie with the same $\frac {3}{28}$ chance, but the third happens $\frac {4}{28}$ of the time, since the unfair coin is heads instead of tails. With 3 possibilities for Jackie and 3 for Phil, there are a total of 9 ways for them both to have 1 head.
            \end{enumerate}
            
            \begin{enumerate}[label=(\roman*)]
                \item Multiplying and adding up all 9 ways, we have a 
 \[\frac {4(3 \cdot 3) + 4(3 \cdot 4) + 1(4 \cdot 4)}{28^{2}} = \frac {100}{784}\]
            \end{enumerate}
            
            \begin{enumerate}[label=(\roman*)]
                \item overall chance for this case.
            \end{enumerate}
            
            \begin{itemize}
        \item \textbf{Case 3}: Two heads.
            \end{itemize}
            
            \begin{enumerate}[label=(\roman*)]
                \item With HHT $\frac {4}{28}$, HTH $\frac {4}{28}$, and THH $\frac {3}{28}$ possible, we proceed as in Case 2, obtaining
            \end{enumerate}
            
 \[\frac {1(3 \cdot 3) + 4(3 \cdot 4) + 4(4 \cdot 4)}{28^{2}} = \frac {121}{784}.\]

            \begin{itemize}
        \item \textbf{Case 4}: Three heads.
            \end{itemize}
            
            \begin{enumerate}[label=(\roman*)]
                \item Similar to Case 1, we can only have HHH, which has $\frac {4}{28}$ chance. Then in this case we get $\frac {16}{784}$
            \end{enumerate}
            Finally, we take the sum: $\frac {9 + 100 + 121 + 16}{784} = \frac {246}{784} = \frac {123}{392}$, so our answer is $123 + 392 = \fbox{515}$.

 


	\item (\textit{AIME I 2010}) Positive integers $a$, $b$, $c$, and $d$ satisfy $a > b > c > d$, $a + b + c + d = 2010$, and $a^2 - b^2 + c^2 - d^2 = 2010$. Find the number of possible values of $a$.


 
 \textbf{Answer: }501



\textbf{Solution 1}

Using the \href{https://artofproblemsolving.com/wiki/index.php?title=Difference\_of\_squares}{difference of squares}, $2010 = (a^2 - b^2) + (c^2 - d^2) = (a + b)(a - b) + (c + d)(c - d) \ge a + b + c + d = 2010$, where equality must hold so $b = a - 1$ and $d = c - 1$. Then we see $a = 1004$ is maximal and $a = 504$ is minimal, so the answer is $\boxed{501}$.

 
Note: We can also find that $b=a-1$ in another way. We know 
 \[a^{2} - b^{2} + c^{2} - d^{2}=(a+b)+(c+d) \implies (a+b)(a-b)-(a+b)+(c+d)(c-d)-(c+d)=0 \implies (a+b)(a-b-1)+(c+d)(c-d-1)=0\]

 Therefore, one of $(a+b)(a-b-1)$ and $(c+d)(c-d-1)$ must be $0,$ because they are both non-negative if $a \geq b + 1$ and $c \geq d + 1$, which is the initial condition. Clearly, $a+b \neq 0$ since then one would be positive and negative, or both would be zero. Therefore, $a-b-1=0$ so $a=b+1$. Similarly, we can deduce that $c=d+1.$




\textbf{Solution 2}

Since $a+b$ must be greater than $1005$, it follows that the only possible value for $a-b$ is $1$ (otherwise the quantity $a^2 - b^2$ would be greater than $2010$). Therefore the only possible ordered pairs for $(a,b)$ are $(504, 503)$, $(505, 504)$, ... , $(1004, 1003)$, so  $a$ has $\boxed{501}$ possible values. Hence the answer is 500+1=501

 


 



\textbf{Solution 3 (Substitution of Differences)}

Let $x = a-b$, $y = b-c$, and $z = c-d$. Since $a > b > c > d \ge 1$, then $x, y, z$ are positive integers.

 Rewrite the second equation using $a^2 - b^2 = (a+b)(a-b)$ and $c^2 - d^2 = (c+d)(c-d)$:

 
 \[x(a+b) + z(c+d) = 2010\]

 Since $a+b+c+d = 2010$, we can substitute $c+d = 2010 - (a+b)$:

 
 \[x(a+b) + z(2010 - (a+b)) = 2010 \implies (x-z)(a+b) = 2010(1-z)\]

 Case 1: $z > 1$

 If $z \ge 2$, then $1-z$ is negative. For the LHS to be negative, we need $z > x$. However, even at minimum values ($x=1, z=2, a+b=1007$), the magnitude of the LHS $|(1-2)1007|$ is far less than $|2010(1-2)|$.

 Case 2: $z = 1$

 If $z = 1$, the RHS becomes $0$. Since $a+b > 0$, we must have $x-z = 0 \implies x = 1$.

 Thus, $a-b = 1$ and $c-d = 1$.

 Finding the Range of $a$:

 Substitute $b=a-1$ and $d=c-1$ into the sum:

 
 \[2a + 2c - 2 = 2010 \implies a + c = 1006\]

 Using the sequence $a > a-1 > c > c-1 \ge 1$:

 $a-1 > c \implies a-1 > 1006 - a \implies 2a > 1007 \implies a \ge 504$
$c-1 \ge 1 \implies 1006 - a \ge 2 \implies a \le 1004$

 Total values = 1004 - 504 + 1 = $\boxed{501}$



	\item (\textit{AIME I 2010}) Let $P(x)$ be a quadratic polynomial with real coefficients satisfying $x^2 - 2x + 2 \le P(x) \le 2x^2 - 4x + 3$ for all real  numbers $x$, and suppose $P(11) = 181$. Find $P(16)$.


 
 \textbf{Answer: }406



\textbf{Solution 1}


                    \begin{figure}[H]
                        \centering
                        \includegraphics[width=0.5\textwidth]{Solution/AIME_I_2010_6_Solution_Figure_1.png}
                    \end{figure}
                    

 Let $Q(x) = x^2 - 2x + 2$, $R(x) = 2x^2 - 4x + 3$. \href{https://artofproblemsolving.com/wiki/index.php?title=Completing\_the\_square}{Completing the square}, we have $Q(x) = (x-1)^2 + 1$, and $R(x) = 2(x-1)^2 + 1$, so it follows that $P(x) \ge Q(x) \ge 1$ for all $x$ (by the \href{https://artofproblemsolving.com/wiki/index.php?title=Trivial\_Inequality}{Trivial Inequality}). 

 Also, $1 = Q(1) \le P(1) \le R(1) = 1$, so $P(1) = 1$, and $P$ obtains its minimum at the point $(1,1)$. Then $P(x)$ must be of the form $c(x-1)^2 + 1$ for some constant $c$; substituting $P(11) = 181$ yields $c = \frac 95$. Finally, $P(16) = \frac 95 \cdot (16 - 1)^2 + 1 = \boxed{406}$.

 



\textbf{Solution 2}

It can be seen that the function $P(x)$ must be in the form $P(x) = ax^2 - 2ax + c$ for some real $a$ and $c$. This is because the \href{https://artofproblemsolving.com/wiki/index.php?title=Derivative}{derivative} of $P(x)$ is $2ax - 2a$, and a global minimum occurs only at $x = 1$ (in addition, because of this derivative, the vertex of any quadratic polynomial occurs at $\frac{-b}{2a}$). Substituting $(1,1)$ and $(11, 181)$ we obtain two equations:

 $P(11) = 99a + c = 181$, and $P(1) = -a + c = 1$.

 Solving, we get $a = \frac{9}{5}$ and $c = \frac{14}{5}$, so $P(x) = \frac{9}{5}x^2 - \frac{18}{5}x + \frac {14}{5}$. Therefore, $P(16) = \boxed{406}$.

 



\textbf{Solution 3}

Let $y = x^2 - 2x + 2$; note that $2y - 1 = 2x^2 - 4x + 3$. Setting $y = 2y - 1$, we find that equality holds when $y = 1$ and therefore when $x^2 - 2x + 2 = 1$; this is true iff $x = 1$, so $P(1) = 1$.

 Let $Q(x) = P(x) - x$; clearly $Q(1) = 0$, so we can write $Q(x) = (x - 1)Q'(x)$, where $Q'(x)$ is some linear function. Plug $Q(x)$ into the given inequality:

 $x^2 - 3x + 2 \le Q(x) \le 2x^2 - 5x + 3$
$(x - 1)(x - 2) \le (x - 1)Q'(x) \le (x - 1)(2x - 3)$, and thus

 $x - 2 \le Q'(x) \le 2x - 3$

 For all $x > 1$; note that the inequality signs are flipped if $x < 1$, and that the division is invalid for $x = 1$. However,

 $\lim_{x \to 1}  x - 2 = \lim_{x \to 1}  2x - 3 = -1$,

 and thus by the \href{https://artofproblemsolving.com/wiki/index.php?title=Sandwich\_theorem}{sandwich theorem} $\lim_{x \to 1} Q'(x) = -1$; by the definition of a continuous function, $Q'(1) = -1$. Also, $Q(11) = 170$, so $Q'(11) = 170/(11-1) = 17$; plugging in and solving, $Q'(x) = (9/5)(x - 1) - 1$. Thus $Q(16) = 390$, and so $P(16) = \boxed{406}$.

 



\textbf{Solution 4}

Let $Q(x) = P(x) - (x^2-2x+2)$, then $0\le Q(x) \le (x-1)^2$ (note this is derived from the given inequality chain).  Therefore, $0\le Q(x+1) \le x^2 \Rightarrow Q(x+1) = Ax^2$ for some real value A.

 $Q(11) = 10^2A \Rightarrow P(11)-(11^2-22+2)=100A \Rightarrow 80=100A \Rightarrow A=\frac{4}{5}$.

 $Q(16)=15^2A=180 \Rightarrow P(16)-(16^2-32+2) = 180 \Rightarrow P(16)=180+226= \boxed{406}$




\textbf{Solution 5}

Let $P(x) = ax^2 + bx + c$. Plugging in $x = 1$ to the expressions on both sides of the inequality, we see that $a + b + c = 1$. We see from the problem statement that $121a + 11b + c = 181$. Since we know the vertex of $P(x)$ lies at $x = 1$, by symmetry we get $81a -9b + c = 181$ as well. Since we now have three equations, we can solve this trivial system and get our answer of $\boxed{406}$.

 



\textbf{Solution 6}

Similar to Solution 5, let $P(x) = ax^2 + bx + c$. Note that $(1,1)$ is a vertex of the polynomial. Additionally, this means that $b = -2a$ (since $\frac{-b}{2a}$ is the minimum $x$ point). Thus, we have $P(x) = ax^2 - 2ax + c$. Therefore $a - 2a + c = 1$. Moreover, $99a + c = 181$. And so our polynomial is $\frac{9}{5}x^2 - \frac{18}{5}x + \frac{14}{5}$. Plug in $x = 16$ to get $\boxed{406}$.

 



\textbf{Solution 7}

Very similar to Solution 6, start by noticing that $P(x)$ is between two polynomials, we try to set them equal to find a point where the two polynomials intersect, meaning that $P(x)$ would also have to intersect that point (it must be between the two graphs). Setting $x^2 - 2x + 2 = 2x^2 - 4x + 3$, we find that $x = 1$. Note that both of these graphs have the same vertex (at $x = 1$), and so $P(x)$ must also have the same vertex $(1, 1)$. Setting $P(x) = ax^2 - 2ax + a + 1$ (this is where we have a vertex at $(1, 1)$), we plug in $11$ and find that $a = 1.8$. Evaluating $1.8x^2 - 3.6x + 2.8$ when $x = 16$ (our intended goal), we find that $P(16) = \boxed{406}$.

 


	\item (\textit{AIME I 2010}) Define an ordered triple $(A, B, C)$ of sets to be \textit{minimally intersecting} if $|A \cap B| = |B \cap C| = |C \cap A| = 1$ and $A \cap B \cap C = \emptyset$. For example, $(\{1,2\},\{2,3\},\{1,3,4\})$ is a minimally intersecting triple. Let $N$ be the number of minimally intersecting ordered triples of sets for which each set is a subset of $\{1,2,3,4,5,6,7\}$. Find the remainder when $N$ is divided by $1000$.


 \textbf{Note}: $|S|$ represents the number of elements in the set $S$.


 
 \textbf{Answer: }760



\textbf{Solution}

Let each pair of two sets have one element in common. Label the common elements as $x$, $y$, $z$. Set $A$ will have elements $x$ and $y$, set $B$ will have $y$ and $z$, and set $C$ will have $x$ and $z$. There are $7 \cdot 6 \cdot 5 = 210$ ways to choose values of $x$, $y$ and $z$. There are $4$ unpicked numbers, and each number can either go in the first set, second set, third set, or none of them. Since we have $4$ choices for each of $4$ numbers, that gives us $4^4 = 256$.

 Finally, $256 \cdot 210 = 53760$, so the answer is $\fbox{760}$.

 \textbf{General Formula} Let this relation be denoted $S_n$ where $n \ge 3$, $S_3 = 6$. Then, we choose the first three elements in $n \cdot (n-1) \cdot (n-2)$ ways. The remaining elements, no element can be in any two of the sets, so either, it is apart of the set $T \not \subset S$, apart of the ordered pair $A$, $B$, or $C$. That is 4 options for all the remaining elements, so $4^{n-3}$ choices. The general formula is $S_n = n \cdot (n-1) \cdot (n-2) \cdot 4^{n-3}$, and plugging in $n=7$, we obtain $S_7 = 7 \cdot 6 \cdot 5 \cdot 4^4 = 256 \cdot 210 = 53\boxed{760}$.

 


	\item (\textit{AIME I 2010}) For a real number $a$, let $\lfloor a \rfloor$ denote the greatest integer less than or equal to $a$. Let $\mathcal{R}$ denote the region in the coordinate plane consisting of points $(x,y)$ such that $\lfloor x \rfloor ^2 + \lfloor y \rfloor ^2 = 25$. The region $\mathcal{R}$ is completely contained in a disk of radius $r$ (a disk is the union of a circle and its interior). The minimum value of $r$ can be written as $\frac {\sqrt {m}}{n}$, where $m$ and $n$ are integers and $m$ is not divisible by the square of any prime. Find $m + n$.


 
 \textbf{Answer: }132



\textbf{Solution}

The desired region consists of 12 boxes, whose lower-left corners are integers solutions of $x^2 + y^2 = 25$, namely $(\pm5,0), (0,\pm5), (\pm3,\pm4), (\pm4,\pm3).$ Since the points themselves are symmetric about $(0,0)$, the boxes are symmetric about $\left(\frac12,\frac12\right)$. The distance from $\left(\frac12,\frac12\right)$ to the furthest point on a box that lays on an axis, for instance $(6,1)$, is $\sqrt {\frac {11}2^2 + \frac12^2} = \sqrt {\frac {122}4}.$ The distance from $\left(\frac12,\frac12\right)$ to the furthest point on a box in the middle of a quadrant, for instance $(5,4)$, is $\sqrt {\frac92^2 + \frac72^2} = \sqrt {\frac {130}4}.$ The latter is the larger, and is $\frac {\sqrt {130}}2$, giving an answer of $130 + 2 = \boxed{132}$.

 
                    \begin{figure}[H]
                        \centering
                        \includegraphics[width=0.5\textwidth]{Solution/AIME_I_2010_8_Solution_Figure_1.png}
                    \end{figure}
                    




\textbf{Solution 2}

When observing the equation $\lfloor x \rfloor ^2 + \lfloor y \rfloor ^2 = 25$, it is easy to see that it is a circle graph. So, we can draw a coordinate plane and find some points. 

 In quadrant $1$, $x < 6$ and $y < 6$. Note that $\lfloor 5.999...\rfloor = 5$, but if we add more $9's$ after the $5$, it will get infinitely close to $6$, so we can use $6$ as a bounding line. Also, with the same logic, when $x = 6$, $y = 1$ (the equal sign represents as $x$ approaches..., not equal to...) So, in quadrant one, we have points $(6,1)$ and $(1,6)$.

 Moving to quadrant $2$, we must note that $\lfloor -4.999...\rfloor = -5$, so the circle will not be centered at $(0,0)$. In quadrant 2, $y$ is still positive, so we can have $y = 1$. When $y = 1$, $x = -5$, so we have our next point $(-5,1)$. With this method, other points can be found in quadrants $3$ and $4$. 

 Additionally, $3 ^2 + 4 ^2 = 5 ^2$, and with the same approaching limit, we know that quadrant $1$ also has lattice points $(4,5)$ and $(5,4)$. We need a point that passes through the circle's center (we don't need to find the center). If we focus on $(5,4)$, the "opposite" point is $(-4,-3)$ located in quadrant 3. 

 Using the distance formula, we find that the distance between the two points is $\sqrt {130}$. Since the line connecting those two points passes through the circle's center, it is the diameter. So, the radius can be found by dividing $\sqrt {130}$ by $2$ to get $\frac {\sqrt {130}}2$, and $m+n=130 + 2 = \boxed{132}$.

 


	\item (\textit{AIME I 2010}) Let $(a,b,c)$ be a real solution of the system of equations $x^3 - xyz = 2$, $y^3 - xyz = 6$, $z^3 - xyz = 20$. The greatest possible value of $a^3 + b^3 + c^3$ can be written in the form $\frac {m}{n}$, where $m$ and $n$ are relatively prime positive integers. Find $m + n$.


 
 \textbf{Answer: }158



\textbf{Solution 1}

Add the three equations to get $a^3 + b^3 + c^3 = 28 + 3abc$.  Now, let $abc = p$.  $a = \sqrt [3]{p + 2}$, $b = \sqrt [3]{p + 6}$ and $c = \sqrt [3]{p + 20}$, so $p = abc = (\sqrt [3]{p + 2})(\sqrt [3]{p + 6})(\sqrt [3]{p + 20})$.  Now \href{https://artofproblemsolving.com/wiki/index.php?title=Cube}{cube} both sides; the $p^3$ terms cancel out.  Solve the remaining \href{https://artofproblemsolving.com/wiki/index.php?title=Quadratic}{quadratic} to get $p = - 4, - \frac {15}{7}$.  To maximize $a^3 + b^3 + c^3$ choose $p = - \frac {15}{7}$ and so the sum is $28 - \frac {45}{7} = \frac {196 - 45}{7}$ giving $151 + 7 = \fbox{158}$.

 



\textbf{Solution 2}

This is almost the same as solution 1. Note $a^3 + b^3 + c^3 = 28 + 3abc$. Next, let $k = a^3$. Note that $b = \sqrt [3]{k + 4}$ and  $c = \sqrt [3]{k + 18}$, so we have $28 + 3\sqrt [3]{k(k+4)(k+18)} = 28+3abc=a^3+b^3+c^3=3k+22$. Move 28 over, divide both sides by 3, then cube to get $k^3-6k^2+12k-8 = k^3+22k^2+18k$. The $k^3$ terms cancel out, so solve the quadratic to get $k = -2, -\frac{1}{7}$. We maximize $abc$ by choosing $k = -\frac{1}{7}$, which gives us $a^3+b^3+c^3 = 3k + 22 = \frac{151}{7}$. Thus, our answer is $151+7=\boxed{158}$.

 



\textbf{Solution 3}

We have that $x^3 = 2 + xyz$, $y^3 = 6 + xyz$, and $z^3 = 20 + xyz$. Multiplying the three equations, and letting $m = xyz$, we have that $m^3 = (2+m)(6+m)(20+m)$, and reducing, that $7m^2 + 43m + 60 = 0$, which has solutions $m = -\frac{15}{7}, -4$. Adding the three equations and testing both solutions, we find the answer of $\frac{151}{7}$, so the desired quantity is $151 + 7 = \fbox{158}$.

 


 


	\item (\textit{AIME I 2010}) Let $N$ be the number of ways to write $2010$ in the form $2010 = a_3 \cdot 10^3 + a_2 \cdot 10^2 + a_1 \cdot 10 + a_0$, where the $a_i$'s are integers, and $0 \le a_i \le 99$. An example of such a representation is $1\cdot 10^3 + 3\cdot 10^2 + 67\cdot 10^1 + 40\cdot 10^0$. Find $N$.


 
 \textbf{Answer: }202



\textbf{Solution 1}

If we choose $a_3$ and $a_1$ such that $(10^3)(a_3) + (10)(a_1) \leq 2010$ there is a \href{https://artofproblemsolving.com/wiki/index.php?title=Bijection}{unique} choice of $a_2$ and $a_0$ that makes the equality hold.  So $N$ is just the number of combinations of $a_3$ and $a_1$ we can pick.  If $a_3 = 0$ or $a_3 = 1$ we can let $a_1$ be anything from $0$ to $99$.  If $a_3 = 2$ then $a_1 = 0$ or $a_1 = 1$.  Thus $N = 100 + 100 + 2 = \fbox{202}$.

 



\textbf{Solution 2}

Note that $a_2\cdot 10^2 + a_0$ is the base $100$ representation of any number from $0$ to $9999$, and similarly $10(a_3\cdot 10^2 + a_1)$ is ten times the base $100$ representation of any number from $0$ to $9999$. Thus, the number of solutions is just the number of solutions to $2010 = 10a+b$ where $0\le a, b\le 9999$, which is equal to $\boxed{202}$ as $a$ can range from $0$ to $201$.

 



\textbf{Solution 3}

Note that $a_0 \equiv 2010\ (\textrm{mod}\ 10)$ and $a_1 \equiv 2010 - a_0\  (\textrm{mod}\ 100)$. It's easy to see that exactly 10 values in $0 \leq a_0 \leq 99$ that satisfy our first congruence. Similarly, there are 10 possible values of $a_1$ for each choice of $a_0$. Thus, there are $10 \times 10 = 100$ possible choices for $a_0$ and $a_1$. We next note that if $a_0$ and $a_1$ are chosen, then a valid value of $a_3$ determines $a_2$, so we dive into some simple casework:

 
            \begin{itemize}
        \item If $2010 - 10a_1 - a_0 \geq 2000$, there are 3 valid choices for $a_3$. There are only 2 possible cases where $2010 - 10a_1 - a_0 \geq 2000$, namely $(a_1, a_0) = (1,0), (10,0)$. Thus, there are $3 \times 2 = 6$ possible representations in this case.
        \item If $2010 - 10a_1 - a_0 < 1000$, $a_3$ can only equal 0. However, this case cannot occur, as $10a_1+a_0\leq 990+99 = 1089$. Thus, $2010-10a_1-a_0 \geq 921$. However,  $2010-10a_1-a_0 = 1000a_3 + 100a_2 \equiv 0\  (\textrm{mod}\ 100)$. Thus, we have $2010-10a_1-a_0 \geq 1000$ always.
        \item If $1000 \leq 2010 - 10a_1 - a_0 < 2000$, then there are 2 valid choices for $a_3$. Since there are 100 possible choices for $a_0$ and $a_1$, and we have already checked the other cases, it follows that $100 - 2 - 0 = 98$ choices of $a_0$ and $a_1$ fall under this case. Thus, there are $2 \times 98 = 196$ possible representations in this case.
            \end{itemize}
            Our answer is thus $6 + 0 + 196 = \boxed{202}$.

 



\textbf{Solution 4: Casework and Brute Force}

We immediately see that $a_3$ can only be $0$, $1$ or $2$. We also note that the maximum possible value for $10a_1 + a_0$ is $990 + 99 = 1089$. We then split into cases:

 Case 1: $a_3 = 0$.We try to find possible values of $a_2$. We plug in $a_3 = 0$ and $10a_1 + a_0 = 1089$ to our initial equation, which gives us $2010 = 0 + 100a_2 + 1089$. Thus $a_2 \geq 10$. We also see that $a_2 \leq 20$. We now take these values of $a_2$ and find the number of pairs $(a_1, a_0)$ that work. If $a_2 = 10$, $10a_1 + a_0 = 1010$. We see that there are $8$ possible pairs in this case. Using the same logic, there are $10$ ways for $a_2 = 11, 12 \ldots 19$. For $a_2 = 20$, we get the equation $10a_1 + a_0 = 10$, for 2 ways. Thus, for $a_3 = 0$, there are $8 + 10 \cdot 9 + 2 = 100$ ways.

 Case 2: $a_3 = 1$.This case is almost identical to the one above, except $0 \leq a_2 \leq 10$. We also get 100 ways.

 Case 3: $a_3 = 2$.If $a_3 = 2$, our initial equation becomes $100a_2 + 10a_1 + a_0 = 10$. It is obvious that $a_2 = 0$, and we are left with $10a_1 + a_0 = 10$. We saw above that there are $2$ ways.

 Totaling everything, we get that there are $100 + 100 + 2 = \boxed{202}$ ways.

 



\textbf{Solution 5: Generating Functions}

We will represent the problem using generating functions. Consider the generating function 
 \[f(x) = (1+x^{1000}+x^{2000}+\cdots+x^{99000})(1+x^{100}+x^{200}+\cdots+x^{9900})(1+x^{10}+x^{20}+\cdots+x^{990})(1+x+x^2+\cdots+x^{99})\]where the first factor represents $a_3$, the second factor $a_2$, and so forth. We want to find the coefficient of $x^{2010}$ in the expansion of $f(x)$. Now rewriting each factor using the geometric series yields 
 \[f(x) = \frac{\cancel{x^{100}-1}}{x-1} \cdot \frac{\cancel{x^{1000}-1}}{x^{10}-1} \cdot \frac{x^{10000}-1}{\cancel{x^{100}-1}} \cdot \frac{x^{100000}-1}{\cancel{x^{1000}-1}}=\frac{x^{10000}-1}{x-1} \cdot \frac{x^{100000}-1}{x^{10}-1} = (1+x+x^2+\cdots + x^{9999})(1+x^{10}+x^{20}+\cdots+x^{99990})\]The coefficient of $x^{2010}$ in this is simply $\boxed{202}$, as we can choose any of the first 202 terms from the second factor and pair it with exactly one term in the first factor.

 



\textbf{Solution 5}

First note that $a_3$ has to be a single-digit number($0$, $1$, or $2$ to be exact), and that $a_1$ has to be a two-digit multiple of ten.Then, $a_3$, $a_2$, $a_1$ and $a_0$ can be represented as follows:
 \begin{align*} a_3 = a \\ a_2 = 10b+c \\ a_1= 10d+e \\ a_0 = 10f \end{align*}
, where $a$, $b$, $c$, $d$, $e$, and $f$ are all(not necessarily nonzero) digits.Now, we can write our given equation as follows:
 \begin{align*} 2010 = 1000(a+b) + 100(c+d) + 10(e+f) \\ 201 = 100(a+b) + 10(c+d) + (e+f) \\ 201 = (100a + 10c + e) + (100b + d + f) \end{align*}
Now, each integer between $0$ and $201$ inclusive can be represented in exactly one way as $100a + 10c + e$, and this corresponds with one unique $100b + d + f$, so it remains to count the number of integers between $0$ and $201$ inclusive. This is easily counted to be $\boxed{202}$.

 



\textbf{Solution 6}

Just note that this corresponds to $0 \leq 10^3 \cdot a_3 + 10^2 \cdot a_2 + 10^1 \cdot a_1 \leq 2010$, because we can use $a_0$ to fill in the remaining gap. Then, dividing by $10$, we have $0 \leq \overline{a_3 a_2 a_1} \leq 201$, of which there are $\boxed{202}$ solutions.

 



\textbf{Solution 7 (Casework, More Expansive, probably the easiest)}

Note that $1000a_3 + 100a_2 + 10a_1$ is always a multiple of 10. Then, if the sum results in 2010, a perfect multiple of 10, that implies $a_0$ must also be a multiple of 10. Therefore, $a_0 = 0, 10, \dots, 90$.

 Imagine we take our $a_0$. Focusing our attention to $1000a_3 + 100a_2 + 10a_1 = 2010 - a_0$, we see we can evenly divide by 10, obtaining $100a_3 + 10a_2 + a_1 = 201-b_0$, where $b_0 = \frac{a_0}{10} = 0, 1, 2, \dots, 9$. Therefore, we must find all ordered triples $(a_3, a_2, a_1)$ such that $100a_3 + 10a_2 + a_1 = 201, 200, \dots, 192$.

 Now, direct our attention to $a_1$. If $100a_3 + 10a_2 + a_1 = 20x, 19y$ where $x = 0, 1$ and $y = 2, 3, \dots, 9$, then $x \equiv x,y \pmod{10}$, dependent on what the value of $100a_3 + 10a_2 + a_1 = 20x, 19y$ is. Now, subtract $a_1$, and divide by 10, and we obtain $10a_3 + a_2$. For all such cases, we either yield $20x - x$ ($x$ is a digit value, not an independent value), which is $200$, and divide by 10, we have 20, OR we yield $19x-x$, which is $190$, and divide by 10, we have 19. There are two cases.

 Case 1: $10a_3 + a_2 = 20$.

 Two such cases will yield this, specifically, $a_0 = 0, 10$. Then, we notice that two subcases arise. Either, $a_1 = 1$, in which we have 20, and so either $a_2 = 0, 10, 20$ and $a_3$ is analogously defined, so three cases here, or $a_1 = 11, 21, 31, \dots, 91$, for which $a_2$ can either be the units digit or the number itself, for which $a_3 = 1,0$, two cases here. There are 9 numbers from $11 - 91$ increment 10, so we have $2 \cdot 9$ cases here, and only one case that yields 3 values, for a total count 21. Two such sets have the property $10a_3 + a_2 = 20$, so we have 42 cases total.

 Case 2: $10a_3 + a_2 = 19$. 

 There are 8 cases that yield this (192 through 199). Then, notice whatever we choose for $a_1$, we will end up with $10a_3 + a_2 = 19, 18, 17, \dots, 10$. Each of these cases have two subcases, either for which $a_2$ can either be the units digit or the number itself, for which $a_3 = 1,0$. There are 10 numbers between $10$ and $19$ inclusive, and that is 20 cases per one. So, per 8, that is 160 cases total.

 Therefore, there are $160 + 42 = \boxed{202}$ ordered triples of integers $(a_3, a_2, a_1, a_0)$ that satisfy the requirement.

 


	\item (\textit{AIME I 2010}) Let $\mathcal{R}$ be the region consisting of the set of points in the coordinate plane that satisfy both $|8 - x| + y \le 10$ and $3y - x \ge 15$. When $\mathcal{R}$ is revolved around the line whose equation is $3y - x = 15$, the volume of the resulting solid is $\frac {m\pi}{n\sqrt {p}}$, where $m$, $n$, and $p$ are positive integers, $m$ and $n$ are relatively prime, and $p$ is not divisible by the square of any prime. Find $m + n + p$.


 
 \textbf{Answer: }365



\textbf{Solution}


                    \begin{figure}[H]
                        \centering
                        \includegraphics[width=0.5\textwidth]{Solution/AIME_I_2010_11_Solution_Figure_1.png}
                    \end{figure}
                    

 The \href{https://artofproblemsolving.com/wiki/index.php?title=Inequality}{inequalities} are equivalent to $y \ge x/3 + 5, y \le 10 - |x - 8|$. We can set them equal to find the two points of intersection, $x/3 + 5 = 10 - |x - 8| \Longrightarrow |x - 8| = 5 - x/3$. This implies that one of $x - 8, 8 - x = 5 - x/3$, from which we find that $(x,y) = \left(\frac 92, \frac {13}2\right), \left(\frac{39}{4}, \frac{33}{4}\right)$. The region $\mathcal{R}$ is a \href{https://artofproblemsolving.com/wiki/index.php?title=Triangle}{triangle}, as shown above. When revolved about the line $y = x/3+5$, the resulting solid is the union of two right \href{https://artofproblemsolving.com/wiki/index.php?title=Cone}{cones} that share the same base and axis. 

 
                    \begin{figure}[H]
                        \centering
                        \includegraphics[width=0.5\textwidth]{Solution/AIME_I_2010_11_Solution_Figure_2.png}
                    \end{figure}
                    

 Let $h_1,h_2$ denote the height of the left and right cones, respectively (so $h_1 > h_2$), and let $r$ denote their common radius. The volume of a cone is given by $\frac 13 Bh$; since both cones share the same base, then the desired volume is $\frac 13 \cdot \pi r^2 \cdot (h_1 + h_2)$. The \href{https://artofproblemsolving.com/wiki/index.php?title=Distance}{distance} from the point $(8,10)$ to the line $x - 3y + 15 = 0$ is given by $\left|\frac{(8) - 3(10) + 15}{\sqrt{1^2 + (-3)^2}}\right| = \frac{7}{\sqrt{10}}$. The distance between $\left(\frac 92, \frac {13}2\right)$ and $\left(\frac{39}{4}, \frac{33}{4}\right)$ is given by $h_1 + h_2 = \sqrt{\left(\frac{18}{4} - \frac{39}{4}\right)^2 + \left(\frac{26}{4} - \frac{33}{4}\right)^2} = \frac{7\sqrt{10}}{4}$. Thus, the answer is $\frac{343\sqrt{10}\pi}{120} = \frac{343\pi}{12\sqrt{10}} \Longrightarrow 343 + 12 + 10 = \boxed{365}$.

 


	\item (\textit{AIME I 2010}) Let $m \ge 3$ be an integer and let $S = \{3,4,5,\ldots,m\}$. Find the smallest value of $m$ such that for every partition of $S$ into two subsets, at least one of the subsets contains integers $a$, $b$, and $c$ (not necessarily distinct) such that $ab = c$.


 \textbf{Note}: a partition of $S$ is a pair of sets $A$, $B$ such that $A \cap B = \emptyset$, $A \cup B = S$.


 
 \textbf{Answer: }243



\textbf{Solution 1}

We claim that $243$ is the minimal value of $m$. Let the two partitioned sets be $A$ and $B$; we will try to partition $3, 9, 27, 81,$ and $243$ such that the $ab=c$ condition is not satisfied. \href{https://artofproblemsolving.com/wiki/index.php?title=Without\_loss\_of\_generality}{Without loss of generality}, we place $3$ in $A$. Then $9$ must be placed in $B$, so $81$ must be placed in $A$, and $27$ must be placed in $B$. Then $243$ cannot be placed in any set, so we know $m$ is less than or equal to $243$.

 For $m \le 242$, we can partition $S$ into $S \cap \{3, 4, 5, 6, 7, 8, 81, 82, 83, 84 ... 242\}$ and $S \cap \{9, 10, 11 ... 80\}$, and in neither set are there values where $ab=c$ (since $8 < (3\text{ to }8)^2 < 81$ and $81^2>242$ and $(9\text{ to }80)^2 > 80$). Thus $m = \boxed{243}$.

 



\textbf{Solution 2}

Consider $\{3,4,12\}$. We could have any two of the three be together in the same set, and the third in the other set. Thus, we have $\{3,4\}, \{3,12\}, \{4,12\}$. We will try to 'place' numbers in either set such that we never have $a\cdot b = c$, until we reach a point where we MUST have $a\cdot b =c$. 

 We begin with $\{3,12\}$. Notice that $a,b,c$ do not have to be distinct, meaning we could have $3\cdot 3=9$. Thus $9$ must be with $4$. Notice that no matter in which set $36$ is placed, we will be forced to have $a\cdot b =c$, since $3*12=36$ and $4*9=36$.

 We could have $\{4,12\}$. Similarly, $16$ must be with $3$, and no matter to which set $48$ is placed into, we will be forced to have $a \cdot b =c$.

 Now we have $\{3,4\}$. $9$ must be with $12$. Then $81$ must be with $\{3,4\}$. Since $27$ can't be placed in the same set as $\{3,4,81\}$, $27$ must go with $\{9,12\}$. But then no matter where $243$ is placed we will have $a\cdot b =c$.

 Thus, $\boxed{243}$ is the minimum $m$.

 


	\item (\textit{AIME I 2010}) Rectangle $ABCD$ and a semicircle with diameter $AB$ are coplanar and have nonoverlapping interiors. Let $\mathcal{R}$ denote the region enclosed by the semicircle and the rectangle. Line $\ell$ meets the semicircle, segment $AB$, and segment $CD$ at distinct points $N$, $U$, and $T$, respectively. Line $\ell$ divides region $\mathcal{R}$ into two regions with areas in the ratio $1: 2$. Suppose that $AU = 84$, $AN = 126$, and $UB = 168$. Then $DA$ can be represented as $m\sqrt {n}$, where $m$ and $n$ are positive integers and $n$ is not divisible by the square of any prime. Find $m + n$.


 
 \textbf{Answer: }069



\textbf{Solution}

\textbf{Diagram} 
                    \begin{figure}[H]
                        \centering
                        \includegraphics[width=0.5\textwidth]{Solution/AIME_I_2010_13_Solution_Figure_1.png}
                    \end{figure}
                    




\textbf{Solution 1}

The center of the semicircle is also the midpoint of $AB$. Let this point be O. Let $h$ be the length of $AD$.

 Rescale everything by 42, so $AU = 2, AN = 3, UB = 4$. Then $AB = 6$ so $OA = OB = 3$.

 Since $ON$ is a radius of the semicircle, $ON = 3$. Thus $OAN$ is an equilateral triangle.

 Let $X$, $Y$, and $Z$ be the areas of triangle $OUN$, sector $ONB$, and trapezoid $UBCT$ respectively.

 $X = \frac {1}{2}(UO)(NO)\sin{O} = \frac {1}{2}(1)(3)\sin{60^\circ} = \frac {3}{4}\sqrt {3}$
$Y = \frac {1}{3}\pi(3)^2 = 3\pi$

 To find $Z$ we have to find the length of $TC$. Project $T$ and $N$ onto $AB$ to get points $T'$ and $N'$. Notice that $UNN'$ and $TUT'$ are similar. Thus:

 $\frac {TT'}{UT'} = \frac {UN'}{NN'} \implies \frac {TT'}{h} = \frac {1/2}{3\sqrt {3}/2} \implies TT' = \frac {\sqrt {3}}{9}h$.

 Then $TC = T'C - T'T = UB - TT' = 4 - \frac {\sqrt {3}}{9}h$. So:

 $Z = \frac {1}{2}(BU + TC)(CB) = \frac {1}{2}\left(8 - \frac {\sqrt {3}}{9}h\right)h = 4h - \frac {\sqrt {3}}{18}h^2$

 Let $L$ be the area of the side of line $l$ containing regions $X, Y, Z$. Then

 $L = X + Y + Z = \frac {3}{4}\sqrt {3} + 3\pi + 4h - \frac {\sqrt {3}}{18}h^2$

 Obviously, the $L$ is greater than the area on the other side of line $l$. This other area is equal to the total area minus $L$. Thus:

 $\frac {2}{1} = \frac {L}{6h + \frac {9}{2}{\pi} - L} \implies 12h + 9\pi = 3L$.

 Now just solve for $h$.

 
 \begin{align*} 12h + 9\pi & = \frac {9}{4}\sqrt {3} + 9\pi + 12h - \frac {\sqrt {3}}{6}h^2 \\ 0 & = \frac {9}{4}\sqrt {3} - \frac {\sqrt {3}}{6}h^2 \\ h^2 & = \frac {9}{4}(6) \\ h & = \frac {3}{2}\sqrt {6} \end{align*}


 Don't forget to un-rescale at the end to get $AD = \frac {3}{2}\sqrt {6} \cdot 42 = 63\sqrt {6}$. 

 Finally, the answer is $63 + 6 = \boxed{069}$.

 



\textbf{Solution 2}

Let $O$ be the center of the semicircle. It follows that $AU + UO = AN = NO = 126$, so triangle $ANO$ is \href{https://artofproblemsolving.com/wiki/index.php?title=Equilateral}{equilateral}.

 Let $Y$ be the foot of the altitude from $N$, such that $NY = 63\sqrt{3}$ and $NU = 21$.

 Finally, denote $DT = a$, and $AD = x$. Extend $U$ to point $Z$ so that $Z$ is on $CD$ and $UZ$ is perpendicular to $CD$. It then follows that $ZT = a-84$. Since $NYU$ and $UZT$ are \href{https://artofproblemsolving.com/wiki/index.php?title=Similar}{similar},

 $\frac {x}{a-84} = \frac {63\sqrt{3}}{21} = 3\sqrt{3}$

 Given that line $NT$ divides $R$ into a ratio of $1:2$, we can also say that

 $(x)(\frac{84+a}{2}) + \frac {126^2\pi}{6} - (63)(21)(\sqrt{3}) = (\frac{1}{3})(252x + \frac{126^2\pi}{2})$

 where the first term is the area of trapezoid $AUTD$, the second and third terms denote the areas of $\frac{1}{6}$ a full circle, and the area of $NUO$, respectively, and the fourth term on the right side of the equation is equal to $R$. Cancelling out the $\frac{126^2\pi}{6}$ on both sides, we obtain

 $(x)(\frac{84+a}{2}) - \frac{252x}{3} = (63)(21)(\sqrt{3})$

 By adding and collecting like terms,$\frac{3ax - 252x}{6} = (63)(21)(\sqrt{3})$
$\frac{(3x)(a-84)}{6} = (63)(21)(\sqrt{3})$.

 Since $a - 84 = \frac{x}{3\sqrt{3}}$,

 $\frac {(3x)(\frac{x}{3\sqrt{3}})}{6} = (63)(21)(\sqrt{3})$
$\frac {x^2}{\sqrt{3}} = (63)(126)(\sqrt{3})$
$x^2 = (63)(126)(3) = (2)(3^5)(7^2)$
$x = AD = (7)(3^2)(\sqrt{6}) = 63\sqrt{6}$, so the answer is $\boxed{069}.$




\textbf{Solution 3}

Note that the total area of $\mathcal{R}$ is $252DA + \frac {126^2 \pi}{2}$ and thus one of the regions has area $84DA + \frac {126^2 \pi}{6}$

 As in the above solutions we discover that $\angle AON = 60^\circ$, thus sector $ANO$ of the semicircle has $\frac{1}{3}$ of the semicircle's area.

 Similarly, dropping the $N'T'$ perpendicular we observe that $[AN'T'D] = 84DA$, which is $\frac{1}{3}$ of the total rectangle.

 Denoting the region to the left of $\overline {NT}$ as $\alpha$ and to the right as $\beta$, it becomes clear that if $[\triangle UT'T] = [\triangle NUO]$ then the regions will have the desired ratio.

 Using the 30-60-90 triangle, the slope of $NT$, is ${-3\sqrt{3}}$, and thus $[\triangle UT'T] = \frac {DA^2}{6\sqrt{3}}$.

 $[NUO]$ is most easily found by $\frac{absin(c)}{2}$: $[\triangle NUO] = \frac {126*42 * \frac {\sqrt{3}}{2}}{2}$

 Equating, $\frac {126*42 * \frac {\sqrt{3}}{2}}{2} = \frac {DA^2}{6\sqrt{3}}$

 Solving, $63 * 21 * 3 * 6 = DA^2$
$DA = 63 \sqrt{6} \longrightarrow \boxed {069}$




\textbf{Solution 4 (Coordinates)}

Like above solutions, note that $ANO$ is equilateral with side length $126,$ where $O$ is the midpoint of $AB.$ Then, if we let $DA=a$ and set origin at $D=(0,0),$ we get $N=(63,a+63\sqrt{3}), U=(84,a).$ Line $NU$ is then $y-a=\sqrt{27}(x-84),$ so it intersects $CA,$ the $x$-axis, at $x=(a/\sqrt{27}+84),$ giving us point $T.$ Now the area of region $R$ is $252a+\pi(126)^2 / 2,$ so one third of that is $84a+\pi(126)^2 / 6.$

 The area of the smaller piece of $R$ is $[AUTD] + [ANU] + [\text{lune} AN] = \frac{1}{2} \cdot a(84+\frac{a}{\sqrt{27}}+84)+\frac{1}{2} \cdot 84 \cdot 63 \sqrt{3}+ \frac{\pi (126)^2}{6}-\frac{1}{2} \cdot 126 \cdot 63 \sqrt{2}$$=\frac{a^2}{2\sqrt{27}}+84a-21\cdot 63\sqrt{2} + \frac{\pi(126)^2}{6}.$Setting this equal to $84a+\pi(126)^2 / 6$ and canceling the $84a + \pi(126)^2$ yields$\frac{a^2}{2\sqrt{27}}=21 \cdot 63 \sqrt{3},$ so $a = 63 \sqrt{6}$ and the anser is $\boxed{069}.$~ rzlng

 



\textbf{Solution 5 (Minimal calculation)}

Once we establish that $\Delta ANO$ is equilateral, we have 
 \[[{\textrm Sector } BON] = 2[{\textrm Sector } AON],  [BCT'U]=2[ADT'U]\]
 \[\Rightarrow [\overset{\large\frown}{NB} CT'UO]=2[\overset{\large\frown}{NA} DT'UO]\]On the other hand,
 \[[\overset{\large\frown}{NB} CT]=2[\overset{\large\frown}{NA} DT]\]Therefore, $[UT'T]=[NUO]$. 

 Now, $UO=42, NU=21 \Rightarrow [UT'T]=[NUO]=2[NN'U]$. Also $\Delta UT'T \sim \Delta NN'U$. Therefore,
 \[DA=UT'=\sqrt{2} NN'=\sqrt{2} \left(\frac{\sqrt{3}}{2}\cdot 126\right)=63\sqrt{6}\longrightarrow \boxed {069}\]



	\item (\textit{AIME I 2010}) For each positive integer $n,$ let $f(n) = \sum_{k = 1}^{100} \lfloor \log_{10} (kn) \rfloor$. Find the largest value of $n$ for which $f(n) \le 300$.


 \textbf{Note:} $\lfloor x \rfloor$ is the greatest integer less than or equal to $x$.


 
 \textbf{Answer: }109



\textbf{Solution 1}

Observe that $f$ is strictly increasing in $n$. We realize that we need $100$ terms to add up to around $300$, so we need some sequence of $2$s, $3$s, and then $4$s.

 It follows that $n \approx 100$ (alternatively, use binary search to get to this, with $n\le 1000$). Manually checking shows that $f(109) = 300$ and $f(110) > 300$. Thus, our answer is $\boxed{109}$.

 



\textbf{Solution 2}

Because we want the value for which $f(n)=300$, the average value of the 100 terms of the sequence should be around $3$. For the value of $\lfloor \log_{10} (kn) \rfloor$ to be $3$, $1000 \le kn < 10000$. We want kn to be around the middle of that range, and for k to be in the middle of 0 and 100, let $k=50$, so $50n=\frac{10000+1000}{2}=\frac{11000}{2}=5500$, and $n = 110$. $f(110) = 301$, so we want to lower $n$. Testing $109$ yields $300$, so our answer is still $\boxed{109}$.

 



\textbf{Solution 3}

For any $n$ where the sum is close to $300$, all the terms in the sum must be equal to $2$, $3$ or $4$. Let $M$ be the number of terms less than or equal to $3$ and $N$ be the number of terms equal to $2$ (also counted in $M$). With this definition of $M$ and $N$ the total will be $400 - M - N \le 300$, from which $M + N \ge 100$. Now $M+1$ is the smallest integer $k$ for which $\log_{10}(kn) \ge 4$ or $kn \ge 10000$, thus 
 \[M = \left\lfloor\frac{9999}{n}\right\rfloor.\] Similarly, 
 \[N = \left\lfloor\frac{999}{n}\right\rfloor = \left\lfloor\frac{M}{10}\right\rfloor.\]

 Therefore, 
 \[M + \left\lfloor \frac{M}{10} \right\rfloor \ge 100 \implies M \ge \left\lceil\frac{1000}{11}\right\rceil = 91 \implies n \le \left\lfloor\frac{9999}{91}\right\rfloor = 109.\] Since we want the largest possible $n$, the answer is  $\boxed{109}$.

 



\textbf{Solution 4 (similar to Solution 1, but with more detail)}

Since we're working with base-$10$ logarithms, we can start by testing out $n$'s that are powers of $10$. For $n = 1$, the terms in the sum are $\lfloor \log_{10} (1)\rfloor, \lfloor \log_{10} (2)\rfloor, \lfloor \log_{10} (3) \rfloor , . . . , \lfloor \log_{10} (100) \rfloor$. For numbers $1$-$9$, $\lfloor \log_{10} (kn) \rfloor = 0$. Then we have $90$ numbers, namely $10$-$99$, for which $\lfloor \log_{10} (kn) \rfloor = 1$. The last number we have is $100$, which gives us $\lfloor \log_{10} (kn) \rfloor = 2$. This sum gives us only $90 + 2 = 92$, which is much too low. However, applying the same counting technique for $n = 10$, our sum comes out to be $9 + 180 + 3 = 192$, since there are $9$ terms for which $\lfloor \log_{10} (kn) \rfloor = 1$, $90$ terms for which $\lfloor \log_{10} (kn) \rfloor = 2$, and one term for which $\lfloor \log_{10} (kn) \rfloor = 3$. So we go up one more power of $10$ and get $18 + 270 + 4 = 292$, which is very close to what we are looking for. 

 Now we only have to bump up the value of $n$ a bit and check our sum. Each increase in $n$ by $1$ actually increases the value of our sum by $1$ as well (except for $n = 101$), because whenever a $4$ is added to the sum, a $3$ is taken away. It doesn't take long to check and see that the value of $n$ we're looking for is $\boxed{109}$ , which corresponds to a sum of exactly $300$.

 


 


	\item (\textit{AIME I 2010}) In $\triangle{ABC}$ with $AB = 12$, $BC = 13$, and $AC = 15$, let $M$ be a point on $\overline{AC}$ such that the incircles of $\triangle{ABM}$ and $\triangle{BCM}$ have equal radii. Then $\frac{AM}{CM} = \frac{p}{q}$, where $p$ and $q$ are relatively prime positive integers. Find $p + q$.


 
 \textbf{Answer: }045



\textbf{Solution 1}

Let $AM = x$, then $CM = 15 - x$. Also let $BM = d$ Clearly, $\frac {[ABM]}{[CBM]} = \frac {x}{15 - x}$. We can also express each area by the rs formula. Then $\frac {[ABM]}{[CBM]} = \frac {p(ABM)}{p(CBM)} = \frac {12 + d + x}{28 + d - x}$. Equating and cross-multiplying yields $25x + 2dx = 15d + 180$ or $d = \frac {25x - 180}{15 - 2x}.$ Note that for $d$ to be positive, we must have $7.2 < x < 7.5$.

 By \href{https://artofproblemsolving.com/wiki/index.php?title=Stewart\%27s\_Theorem}{Stewart's Theorem}, we have $12^2(15 - x) + 13^2x = d^215 + 15x(15 - x)$ or $432 = 3d^2 + 40x - 3x^2.$ Brute forcing by plugging in our previous result for $d$, we have $432 = \frac {3(25x - 180)^2}{(15 - 2x)^2} + 40x - 3x^2.$ Clearing the fraction and gathering like terms, we get $0 = 12x^4 - 340x^3 + 2928x^2 - 7920x.$
<i>Aside: Since $x$ must be rational in order for our answer to be in the desired form, we can use the \href{https://artofproblemsolving.com/wiki/index.php?title=Rational\_Root\_Theorem}{Rational Root Theorem} to reveal that $6x$ is an integer because we can divide the polynomial by $2$. The only such $x$ in the above-stated range is $\frac {22}3$.</i>

 Legitimately solving that quartic, note that $x = 0$ and $x = 15$ should clearly be solutions, corresponding to the sides of the triangle and thus degenerate cevians. Factoring those out, we get $0 = 4x(x - 15)(3x^2 - 40x + 132) = x(x - 15)(x - 6)(3x - 22).$ The only solution in the desired range is thus $\frac {22}3$. Then $CM = \frac {23}3$, and our desired ratio $\frac {AM}{CM} = \frac {22}{23}$, giving us an answer of $\boxed{045}$.

 



\textbf{Solution 2}

Let $AM = 2x$ and $BM = 2y$ so $CM = 15 - 2x$. Let the \href{https://artofproblemsolving.com/wiki/index.php?title=Incenter}{incenters} of $\triangle ABM$ and $\triangle BCM$ be $I_1$ and $I_2$ respectively, and their equal inradii be $r$. From $r = \sqrt {(s - a)(s - b)(s - c)}/s$, we find that

 
 \begin{align*}r^2 & = \frac {(x + y - 6)( - x + y + 6)(x - y + 6)}{x + y + 6} & (1) \\ & = \frac {( - x + y + 1)(x + y - 1)( - x - y + 14)}{ - x + y + 14}. & (2) \end{align*}


 Let the incircle of $\triangle ABM$ meet $AM$ at $P$ and the incircle of $\triangle BCM$ meet $CM$ at $Q$. Then note that $I_1 P Q I_2$ is a rectangle. Also, $\angle I_1 M I_2$ is right because $MI_1$ and $MI_2$ are the angle bisectors of $\angle AMB$ and $\angle CMB$ respectively and $\angle AMB + \angle CMB = 180^\circ$. By properties of \href{https://artofproblemsolving.com/wiki/index.php?title=Tangent\_(geometry)}{tangents} to \href{https://artofproblemsolving.com/wiki/index.php?title=Circle}{circles} $MP = (MA + MB - AB)/2 = x + y - 6$ and $MQ = (MB + MC - BC)/2 = - x + y + 1$. Now notice that the altitude of $M$ to $I_1 I_2$ is of length $r$, so by similar triangles we find that $r^2 = MP \cdot MQ = (x + y - 6)( - x + y + 1)$ (3). Equating (3) with (1) and (2) separately yields

 
 \begin{align*} 2y^2 - 30 = 2xy + 5x - 7y \\  2y^2 - 70 = - 2xy - 5x + 7y, \end{align*}


 and adding these we have

 
 \[4y^2 - 100 = 0\implies y = 5\implies x = 11/3 \\ \implies AM/MC = (22/3)/(15 - 22/3) = 22/23 \implies \boxed{045}.\]




\textbf{Solution 3}

Let the incircle of $ABM$ hit $AM$, $AB$, $BM$ at $X_{1},Y_{1},Z_{1}$, and let the incircle of $CBM$ hit $MC$, $BC$, $BM$ at $X_{2},Y_{2},Z_{2}$. Draw the incircle of $ABC$, and let it be tangent to $AC$ at $X$. Observe that we have a homothety centered at A sending the incircle of $ABM$ to that of $ABC$, and one centered at $C$ taking the incircle of $BCM$ to that of $ABC$. These have the same power, since they take incircles of the same size to the same circle. Also, the power of the homothety is $AX_{1}/AX=CX_{2}/CX$.

 By standard computations, $AX=\dfrac{AB+AC-BC}{2}=7$ and $CX=\dfrac{BC+AC-AB}{2}=8$. Now, let $AX_{1}=7x$ and $CX_{2}=8x$. We will now go around and chase lengths. Observe that $BY_{1}=BA-AY_{1}=BA-AX_{1}=12-7x$. Then, $BZ_{1}=12-7x$. We also have $CY_{2}=CX_{2}=8x$, so $BY_{2}=13-8x$ and $BZ_{2}=13-8x$.

 Observe now that $X_{1}M+MX_{2}=AC-15x=15(1-x)$. Also,$X_{1}M-MX_{2}=MZ_{1}-MZ_{2}=BZ_{2}-BZ_{1}=BY_{2}-BY_{1}=(1-x)$. Solving, we get $X_{1}M=8-8x$ and $MX_{2}=7-7x$ (as a side note, note that $AX_{1}+MX_{2}=X_{1}M+X_{2}C$, a result that I actually believe appears in Mandelbrot 1995-2003, or some book in that time-range).

 Now, we get $BM=BZ_{2}+Z_{2}M=BZ_{2}+MX_{2}=20-15x$. To finish, we will compute area ratios. $\dfrac{[ABM]}{[CBM]}=\dfrac{AM}{MC}=\dfrac{8-x}{7+x}$. Also, since their inradii are equal, we get $\dfrac{[ABM]}{[CBM]}=\dfrac{40-16x}{40-14x}$. Equating and cross multiplying yields the quadratic $3x^{2}-8x+4=0$, so $x=2/3,2$. However, observe that $AX_{1}+CX_{2}=15x<15$, so we take $x=2/3$. Our ratio is therefore $\dfrac{8-2/3}{7+2/3}=\dfrac{22}{23}$, giving the answer $\boxed{045}$.

 Note: Once we have $MX_1=8-8x$ and $MX_2=7-7x$, it's bit easier to use use the right triangle of $O_1MO_2$ than chasing the area ratio. The inradius of $\triangle{ABC}$ can be calculated to be $r=\sqrt{14}$, and the inradius of $ABM$ and $ACM$ are $r_1=r_2= xr$, so, 
 \[O_1O_2^2 = O_1M^2 + O_2M^2 = r_1^2+X_1M^2 + r_2^2 + X_2M^2\]or,
 \[(15(1-x))^2 = 2(\sqrt{14}x)^2 + (7(1-x))^2 + (8(1-x))^2\]
 \[112(1-x)^2 = 28x^2\]
 \[4(1-x)^2 = x^2\]We get $x=\frac{2}{3}$ or $x=2$.

 



\textbf{Solution 4}

Suppose the incircle of $ABM$ touches $AM$ at $X$, and the incircle of $CBM$ touches $CM$ at $Y$. Then

 
 \[r = AX \tan(A/2) = CY \tan(C/2)\]

 We have $\cos A = \frac{12^2+15^2-13^2}{2\cdot 12\cdot 15} = \frac{200}{30\cdot 12}=\frac{5}{9}$, $\tan(A/2) = \sqrt{\frac{1-\cos A}{1+\cos A}} = \sqrt{\frac{9-5}{9+5}} = \frac{2}{\sqrt{14}}$
$\cos C = \frac{13^2+15^2-12^2}{2\cdot 13\cdot 15} = \frac{250}{30\cdot 13} = \frac{25}{39}$, $\tan(C/2) = \sqrt{\frac{39-25}{39+25}}=\frac{\sqrt{14}}{8}$,

 Therefore $AX/CY = \tan(C/2)/\tan(A/2) = \frac{14}{2\cdot 8}= \frac{7}{8}.$

 And since $AX=\frac{1}{2}(12+AM-BM)$, $CY = \frac{1}{2}(13+CM-BM)$, 

 
 \[\frac{12+AM-BM}{13+CM-BM} = \frac{7}{8}\]

 \[96+8AM-8BM = 91 +7CM-7BM\]

 \[BM= 5 + 8AM-7CM = 5 + 15AM - 7(CM+AM) = 5+15(AM-7) \dots\dots (1)\]

 Now,

 $\frac{AM}{CM} = \frac{[ABM]}{[CBM]} = \frac{\frac{1}{2}(12+AM+BM)r}{\frac{1}{2}(13+CM+BM)r}=\frac{12+AM+BM}{13+CM+BM}= \frac{12+BM}{13+BM} = \frac{17+15(AM-7)}{18+15(AM-7)}$

 \[\frac{AM}{15} = \frac{17+15(AM-7)}{35+30(AM-7)} = \frac{15AM-88}{30AM-175}\]
 \[6AM^2 - 35AM = 45AM-264\]
 \[3AM^2 -40AM+132=0\]
 \[(3AM-22)(AM-6)=0\]

 So $AM=22/3$ or $6$. But from (1) we know that $5+15(AM-7)>0$, or $AM>7-1/3>6$, so $AM=22/3$, $CM=15-22/3=23/3$, $AM/CM=22/23$.

 



\textbf{Solution 5}

Let the common inradius equal r, $BM = x$, $AM = y$, $MC = z$

 From the prespective of $\triangle{ABM}$ and $\triangle{BMC}$ we get:

 $S_{ABM} = rs = r \cdot (\frac{12+x+y}{2})$ $S_{BMC} = rs = r \cdot (\frac{13+x+z}{2})$

 Add two triangles up, we get $\triangle{ABC}$ :

 $S_{ABC} =  S_{ABM} +  S_{BMC} = r \cdot \frac{25+2x+y+z}{2}$

 Since $y + z = 15$, we get:

 \begin{verbatim}
$r = \frac{S_{ABC}}{20 + x}$
\end{verbatim}

By drawing an altitude from $I_1$ down to a point $H_1$ and from $I_2$ to $H_2$, we can get:

 $r \cdot cot(\frac{\angle A}{2}) =r \cdot A H_1 = r \cdot \frac{12+y-x}{2}$  and

 $r \cdot cot(\frac{\angle C}{2}) = r \cdot H_2 C =  r \cdot \frac{13+z-x}{2}$

 Adding these up, we get:

 $r \cdot (cot(\frac{\angle A}{2})+cot(\frac{\angle C}{2})) = \frac{25+y+z-2x}{2} = \frac{40-2x}{2} = 20-x$
\begin{verbatim}
$r = \frac{20-x}{cot(\frac{\angle A}{2})+cot(\frac{\angle C}{2})}$
\end{verbatim}

Now, we have 2 values equal to r, we can set them equal to each other:

 \begin{verbatim}
$\frac{S_{ABC}}{20 + x} = \frac{20-x}{cot(\frac{\angle A}{2})+cot(\frac{\angle C}{2})}$ 
\end{verbatim}

If we let R denote the incircle of ABC, note:

 AC = $(cot(\frac{\angle A}{2})+cot(\frac{\angle C}{2})) \cdot R = 15$ and

 $S_{ABC} = \frac{12+13+15}{2} \cdot R = 20 \cdot R$. 

 By cross multiplying the equation above, we get:

 \begin{verbatim}
$400 - x^2 = (cot(\frac{\angle A}{2})+cot(\frac{\angle C}{2})) \cdot S_{ABC} = (cot(\frac{\angle A}{2})+cot(\frac{\angle C}{2})) \cdot R \cdot 20 = 15 \cdot 20 = 300$
\end{verbatim}

We can find out x:

 \begin{verbatim}
$x = 10$.
\end{verbatim}

Now, we can find ratio of y and z:

 \begin{verbatim}
$\frac{AM}{CM} = \frac{y}{z} = \frac{S_{ABM}}{S_{BCM}} = \frac{r \cdot \frac{22+y}{2} }{r \cdot \frac{23+z}{2}} = \frac{22+y}{23+z} = \frac{22}{23}$ 
\end{verbatim}

The answer is $\boxed{045}$. 

 -Alexlikemath

 



\textbf{Solution 6 (Similar to Solution 1 with easier computation)}

Let $CM=x, AM=rx, BM=d$. $x+rx=15\Rightarrow x=\frac{15}{1+r}$.

 Similar to Solution 1, we have 
 \[r=\frac{[AMB]}{[CMB]}=\frac{12+rx+d}{13+x+d} \Rightarrow d=\frac{13r-12}{1-r}\]as well as
 \[12^2\cdot x + 13^2 rx=15x\cdot rx+15d^2 (\text{via Stewart's Theorem})\]
 \[\frac{(12^2 + 13^2r) \cdot 15}{1+r} - \frac{15r\cdot 15^2}{(1+r)^2}=\frac{15(13r-12)^2}{(1-r)^2}\]
 \[\frac{169r^2+88r+144}{(1+r)^2}=\frac{(13r-12)^2}{(1-r)^2} =\frac{169r^2-312r+144}{(1-r)^2} =\frac{400r}{4r}=100\](here we used the fact that if $\frac{a}{b} = \frac{c}{d} = k,$ then $\frac{a-c}{b-d}=k$ as well.)

 Notice $\frac{12}{13} < r < 1$, so $\frac{13r-12}{1-r} = 10$ and $\boxed{r = \frac{22}{23}}.$




\textbf{Solution 7 (No Stewart's or trig, fast + clever)}

Let $BM = d, AM = x, CM = 15 - x$. Observe that we have the equation by the incircle formula:
 \[\frac{[ABM]}{12 + AM + MB} = \frac{[CBM]}{13 + CM + MB} \implies \frac{AM}{CM} = \frac{12 + MB}{13 + MB} \implies \frac{x}{15 - x} = \frac{12 + d}{13 + d}.\]Now let $X$ be the point of tangency between the incircle of $\triangle ABC$ and $AC$. Additionally, let $P$ and $Q$ be the points of tangency between the incircles of $\triangle ABM$ and $\triangle CBM$ with $AC$ respectively. Some easy calculation yields $AX = 7, CX = 8$. By homothety we have
 \[\frac{AP}{7} = \frac{CQ}{8} \implies 8(AP) = 7(CQ) \implies 8(12 + x - d) = 7(13 + 15 - x - d) \implies d = 15x - 100.\]Substituting into the first equation derived earlier it is left to solve
 \[\frac{x}{15 - x} = \frac{15x - 88}{15x - 87} \implies 3x^2 - 40x + 132 \implies (x - 6)(3x - 22) = 0.\]Now $x = 6$ yields $d = -10$ which is invalid, hence $x = \frac{22}{3}$ so $\frac{AM}{CM} = \frac{\frac{22}{3}}{15 - \frac{22}{3}} = \frac{22}{23}.$ The requested sum is $22 + 23 = \boxed{45}$. ~blueprimes

 


\end{enumerate}